%% ma: L10-B13-0093
%% lop: 10
%% chuong: 5
%% bai: 13
%% dang: 10.13.11
%% loai: TN4
%% muc: TH
%% dap_an: "D"
%% nguon: "chuyen-de-cac-so-dac-trung-cua-mau-so-lieu-khong-ghep-nhom-toan-10-chuong-trinh-moi.pdf (GV. Phan Nhật Linh) (P03), Bài 02 Dạng 2 Phần I Câu 20"
%% kien_thuc: "Lớp 10 Bài 13 (trung vị, tứ phân vị)"
%% trang_thai: chinh-thuc
%% ghi_chu: "Tự duyệt theo yêu cầu giáo viên 22/10/2026."
\begin{ex}
Số nhân khẩu của các hộ gia đình ở một khu vực được thống kê như sau: $2$; $5$; $3$; $7$; $11$; $12$; $4$; $6$; $5$. Tứ phân vị thứ ba của mẫu số liệu là
\choice{$10$}{$11$}{$4$}{\True $9$}
\loigiai{Sắp xếp: $2$; $3$; $4$; $5$; $5$; $6$; $7$; $11$; $12$; $Q_2=5$. Nửa trên $6;7;11;12$ nên $Q_3=\dfrac{7+11}2=9$. Chọn D.}
\end{ex}
